\documentclass[11pt]{article}
\usepackage[margin=1in]{geometry}
\usepackage{amsmath,amssymb,amsthm,mathtools}
\usepackage[hidelinks]{hyperref}
\newtheorem{theorem}{Theorem}
\newtheorem{corollary}[theorem]{Corollary}
\newtheorem{lemma}[theorem]{Lemma}
\theoremstyle{remark}\newtheorem{remark}[theorem]{Remark}
\newcommand{\Ric}{\operatorname{Ric}}
\newcommand{\Vol}{\operatorname{Vol}}
\newcommand{\Var}{\operatorname{Var}}
\newcommand{\Yam}{\mathcal Y}
\title{Quantitative Scalar-Curvature Excursions from an All-Metric Ricci-Defect Gap}
\author{Qiushi Line 1: external-theorem application note}
\date{7 October 2026}
\begin{document}
\maketitle

\begin{abstract}
Assuming the four-dimensional scalar-curvature extension theorem stated
in the specified OpenAI manuscript, an all-metric positive lower bound
for the squared trace-free Ricci curvature on a negative-Yamabe
four-manifold forces a volume-normalized Ricci flow to reach every
positive scalar-curvature threshold. We give an explicit upper bound
for the first hitting time. The elementary integral estimate also gives
linear growth of the running scalar-curvature maximum for any immortal
flow. This is a downstream consequence of the existing Line 1 static
theorems, not a solution of the unrestricted O10 endpoint or the O11
optimal leading-coefficient problem. No priority claim is made.
\end{abstract}

\section{Inputs and conventions}
Let $M$ be a closed connected smooth four-manifold. Assume
\begin{equation}
 P(g):=\int_M|\Ric_g^\circ|^2\,dV_g\ge p_0>0
 \quad\hbox{for every smooth metric }g,
 \qquad \Yam(M)<0.\label{eq:static}
\end{equation}
The all-metric hypothesis is essential; a positive lower bound for
$P+D_2$ alone is not substituted for it.

We use the following external input, referred to as (F): a smooth
unnormalized Ricci flow on a closed four-manifold extends smoothly
past a finite time if its scalar curvature remains uniformly bounded
in absolute value. This is Theorem~\texttt{thm:extension} in
\cite{OAIFlow}, at the fixed source revision specified below.
Our application takes (F) as an external theorem. It does not constitute
an independent verification of its full proof.

Let $g(t)$ be the maximal volume-normalized flow
\begin{equation}
 \partial_tg=-2\Ric+\tfrac12\bar s\,g,
 \qquad \bar s(t)=V^{-1}\int_Ms_{g(t)}\,dV_{g(t)},
 \qquad V=\Vol_{g(0)}M.\label{eq:flow}
\end{equation}
The volume $V$ is constant. Put
\[
 m(t)=\min_Ms_{g(t)},\qquad m_0=m(0),\qquad
 b_0=-m_0,\qquad\bar s_0=\bar s(0).
\]
Choose $a>0$ such that $\Yam(M)\le-a\sqrt V$.
The Laplacian on functions is $\Delta=\operatorname{div}\nabla$.

\section{A finite-time integral estimate}
\begin{lemma}\label{lem:min}
Throughout the smooth flow, $m_0\le m(t)\le-a<0$. At almost every time,
\begin{equation}
 m'(t)\ge\frac{-m(t)}2\bigl(\bar s(t)-m(t)\bigr)\ge0.
 \label{eq:min}
\end{equation}
\end{lemma}
\begin{proof}
For a metric with scalar minimum $m<0$, the Yamabe quotient satisfies
\[
 \frac{\int(6|du|^2+s u^2)}{(\int u^4)^{1/2}}
 \ge m\frac{\int u^2}{(\int u^4)^{1/2}}
 \ge m\sqrt V.
\]
Here the second inequality uses the negative sign of $m$ and Holder's
inequality. Thus $m\sqrt V\le Y([g])\le\Yam(M)\le-a\sqrt V$.
A nonnegative scalar minimum would instead imply a nonnegative Yamabe
constant, contradicting $\Yam(M)<0$.

The scalar evolution for \eqref{eq:flow} is
\begin{equation}
 \partial_ts=\Delta s+2|\Ric^\circ|^2+\tfrac12s(s-\bar s).
 \label{eq:scalar}
\end{equation}
On every compact smooth time interval, $m(t)$ is locally Lipschitz;
the minimum principle gives \eqref{eq:min} almost everywhere.
It follows that $m(t)\ge m_0$. See \cite[Sections 2--3]{Ishida}
for the underlying scalar-minimum argument.
\end{proof}

\begin{lemma}\label{lem:integral}
If $K\ge0$ and $s_{g(u)}\le K$ for $0\le u\le t$, then
\begin{equation}
 \frac{2p_0}{V}t\le
 K\left(1+\log\frac{b_0}{a}\right)-\bar s_0+b_0-a.
 \label{eq:window}
\end{equation}
This conclusion does not use the extension theorem (F).
\end{lemma}
\begin{proof}
Since $\partial_t dV=(\bar s-s)dV$, integration of
\eqref{eq:scalar} gives the exact identity
\begin{equation}
 \bar s'=\frac{2P(g(t))}{V}-\frac12\Var(s),
 \qquad \Var(s)=V^{-1}\int(s-\bar s)^2.
 \label{eq:mean}
\end{equation}
The pointwise inequality $(s-m)(K-s)\ge0$ implies
\[
 \Var(s)\le(K-\bar s)(\bar s-m)
             \le(K-m)(\bar s-m).
\]
By Lemma~\ref{lem:min},
\[
 \frac12\Var(s)\le\left(1+\frac{K}{-m}\right)m'.
\]
Write $u_t=-m(t)$. Integration yields
\begin{align}
 \frac12\int_0^t\Var(s)\,du
 &\le K\log\frac{b_0}{u_t}+b_0-u_t\notag\\
 &\le K\log\frac{b_0}{a}+b_0-a,\label{eq:variance-integral}
\end{align}
because $u_t\ge a$ and $K\ge0$.
Integrate \eqref{eq:mean}, use $P\ge p_0$ and $\bar s(t)\le K$,
and apply \eqref{eq:variance-integral}. This proves
\eqref{eq:window}.
\end{proof}

\section{Thresholds and running maxima}
\begin{theorem}[Quantitative threshold crossing]\label{thm:hitting}
Assume \eqref{eq:static} and the external theorem (F). For every
$K>\max\{0,\max_Ms_{g(0)}\}$, the first hitting time
\[
 \tau_K=\inf\{t:\max_Ms_{g(t)}\ge K\}
\]
exists strictly before the maximal smooth time and satisfies
\begin{equation}
 \boxed{\displaystyle
 \tau_K\le\frac{V}{2p_0}
 \left[K\left(1+\log\frac{b_0}{a}\right)-\bar s_0+b_0-a\right].}
 \label{eq:hitting}
\end{equation}
\end{theorem}
\begin{proof}
Suppose the flow has a finite maximal time $T$ before it reaches $K$.
Lemma~\ref{lem:min} bounds its scalar curvature between $m_0$ and $K$.
Set
\[
 c(t)=\exp\left(-\tfrac12\int_0^t\bar s(u)\,du\right),
 \qquad \widetilde t=\int_0^t c(u)\,du,
 \qquad\widetilde g(\widetilde t)=c(t)g(t).
\]
The factor $c$ has positive upper and lower bounds on $[0,T)$.
Direct differentiation gives
$\partial_{\widetilde t}\widetilde g=-2\Ric_{\widetilde g}$.
The transformed terminal time is finite and
$s_{\widetilde g}=s_g/c$ is bounded. By (F), the transformed flow
extends smoothly. Restoring its constant volume and time parameter
extends $g$, a contradiction.

If the threshold is never reached and the flow is immortal,
Lemma~\ref{lem:integral} contradicts $t\to\infty$.
Thus $\tau_K$ exists before the maximal time. Apply
\eqref{eq:window} for $t<\tau_K$ and pass to $\tau_K$.
\end{proof}

\begin{corollary}[Running scalar maximum]
Define
\[
 K(t)=\max\{0,\sup_{0\le u\le t}\max_Ms_{g(u)}\}.
\]
Throughout smooth existence,
\begin{equation}
 \boxed{\displaystyle
 K(t)\ge\frac{(2p_0/V)t+\bar s_0-b_0+a}
 {1+\log(b_0/a)}.}\label{eq:growth}
\end{equation}
In particular, if the flow is immortal, its running scalar maximum
has a positive linear lower growth rate. Formula~\eqref{eq:growth}
is not a finite-time blowup rate and does not assert monotonicity of
the scalar maximum at individual times.
\end{corollary}
\begin{proof}
Apply Lemma~\ref{lem:integral} with the interval bound $K(t)$ and
rearrange.
\end{proof}

\section{Insertion into the Line 1 hypotheses}
Suppose the actual supply and the Line 1 estimates yield, for every
metric,
\[
 \mathcal J(g)=4P(g)/R(g)\ge j_0>0,
 \qquad S_2(g)\ge32\pi^2 A,
 \qquad\Yam(M)\le-4\pi\sqrt{2A}.
\]
Then
\[
 R=P+S_2/4\ge8\pi^2A,\qquad
 p_0=2\pi^2Aj_0,\qquad a=4\pi\sqrt{2A/V}
\]
are admissible. In an N1 application, one must first ensure that
its topological excess leaves a strictly positive $j_0$; an estimate
for $\mathcal J+d$ is not automatically a positive estimate for
$\mathcal J$ alone.

If $s_{g(0)}=-4$, then
\begin{equation}
 \tau_K\le\frac{V}{2p_0}
 \left[K\left(1+\log\frac4a\right)+8-a\right].
 \label{eq:constant-initial}
\end{equation}
On the original balanced family, use the proved linear lower bound
$p_{0,j}\ge cA_jh_j$. For the selected scalar-minus-four initial
representatives, a uniform bound $V_j\le C A_j$ then gives
\[
 \tau_{K,j}\le C_K/h_j.
\]
The volume bound follows in the manuscript's low-defect window from
$8\pi^2A+P\le R\le8\pi^2AU+4P$, $U=1+O(h)$,
$\mathcal J=O(h)$, and $V=(R-P)/4$.
This is only an upper bound, not an asymptotically sharp lifetime.
It uses the established linear scale, not the unknown coefficient in O11.

\begin{remark}[Scope and novelty]
The topology--gauge obstruction to long-time scalar-bounded normalized
Ricci flow predates this note; see \cite{Ishida}. The new source used
here removes the possible finite-time scalar-bounded termination.
The displayed quantitative threshold estimate is derived explicitly
above. No claim of global bibliographic novelty, full certification of
the external manuscript, unrestricted O10 closure, or O11 sharpness
is made.
\end{remark}

\section{Two precise cautions for the remaining transfers}
The tamed-to-compatible theorem in \cite{OAITame} starts with a globally
smooth almost-complex structure and an everywhere taming closed form.
It does not supply a curvature-energy comparison for the new compatible
form. In the original surgery family,
\[
 a_\epsilon^2=A-k,\quad T_+=A-k-4\ell,\quad
 c_2(S^+_{a_\epsilon})[M]=\frac{a_\epsilon^2-T_+}{4}=\ell.
\]
For $\ell>0$ a nowhere-zero section of this particular positive spinor
bundle is impossible. Thus one cannot normalize its actual spinor
globally merely to invoke that theorem. This is not a claim that no
other almost-complex structure exists on $M$. The density-splitting
mechanism may still be relevant, but the required actual current,
closure, geometric bounds and useful cohomological pairing remain to
be established in the degenerating O10 setting.

The tree theorem in \cite{OAITree} fixes a finite Ricci-flat tree,
nontrivial matching quotient groups and high-order weighted errors.
Its renormalized Einstein--Hilbert functional has length weight two;
$P$ is scale invariant in dimension four. Its constants need not be
uniform over growing trees, and a scalar-flat surgery core need not be
Ricci-flat. A transfer to O11 therefore needs a comparison for the
correct functional, a zero-order topological reference value and
summable errors $o(Q)$; near-stationarity alone does not supply them.

\begin{thebibliography}{9}
\bibitem{OAIFlow}
OpenAI, \emph{Bounded scalar curvature and smooth extension of
four-dimensional Ricci flow}, September 24, 2026.
Repository revision \texttt{adc7f1241b42e322a6451854ab7e4b4c146bf78a};
\texttt{01-introduction.tex}, Theorem \texttt{thm:extension}.
\url{https://github.com/openai/math}

\bibitem{Ishida}
M. Ishida, \emph{The Normalized Ricci Flow on Four-Manifolds and
Exotic Smooth Structures}, arXiv:0807.2169, Sections 2--3.
\url{https://arxiv.org/abs/0807.2169}

\bibitem{OAITame}
OpenAI, \emph{Taming implies compatibility on four-manifolds},
September 23, 2026. Same repository revision;
\texttt{01-introduction.tex}, \texttt{04-energy.tex}, and
\texttt{06-splitting.tex}.

\bibitem{OAITree}
OpenAI, \emph{Path selection and an elliptic inequality on degenerating
Ricci-flat trees}, September 24, 2026. Same repository revision;
\texttt{02-tree.tex} and \texttt{07-observables.tex}.

\bibitem{Qiushi}
Qiushi Line 1 project manuscript and inheritance notes:
\emph{20260929 Line1 Integrated Manuscript},
\texttt{eq:vg-two-sided-R}, the Spin$^c$ characteristic-number identity,
and \emph{03 N1 Complete Inheritance}. Used as project premises,
not re-certified in this note.
\end{thebibliography}
\end{document}
